% ============================================================
% Section 7: Related Work
% ============================================================
\section{Related Work}
\label{sec:related_work}

\widefloatbegin
\centering
\caption{Comparison of representative modeling attack methods on XOR APUFs.
CRP budgets are approximate and reported at or near the transition midpoint.}
\label{tab:related_comparison}
\resizebox{\textwidth}{\height}{%
\begin{tabular}{lcccccl}
\toprule
Method & Year & $k$-range & CRP required & Model & Notes \\
\midrule
LR (Ruhrmair~et~al.)~\cite{ruhrmair2010modeling}
  & 2010 & 1--6 & $10^4$--$2\times10^5$ & LR & parity-encoded features \\
CMA-ES (Becker)~\cite{becker2015cia}
  & 2015 & 1--32 & $4\times10^3$--$2\times10^6$ & CMA-ES & reliability-based \\
Split Attack (Wisiol~et~al.)~\cite{wisiol2020splitting}
  & 2020 & i (iPUF) & $\sim$$10^6$ & MLP & Breaks iPUF \\
MoPE (Fei~et~al.)~\cite{fei2024mope}
  & 2024 & 1--7 & $10^4$--$2.4\times10^6$ & MLP+MoPE & Improves low-XOR \\
Chosen Challenge (Sayadi~et~al.)~\cite{sayadi2025}
  & 2025 & 1--64 & $10^2$--$10^5$ & LR & chosen challenges \\
CalyPSO (Wisiol~et~al.)~\cite{calypso2024}
  & 2024 & 1--20 & $10^5$--$10^7$ & PSO & GPU-accelerated \\
This work & 2026 & 2--10 & $10^5$--$10^7$ & MLP & Probabilistic S-curve \\
\bottomrule
\end{tabular}
}
\widefloatend

\paragraph{Traditional ML attacks.}
Modeling attacks on Arbiter PUFs were first demonstrated by
Ruhrmair~et~al.~\cite{ruhrmair2010modeling}, who showed that logistic regression
with parity-encoded features achieves high accuracy on single APUFs, establishing
the linear separability of individual APUF responses and motivating the XOR APUF
as a countermeasure.  Subsequent work attacked higher XOR configurations through
evolution strategies: Becker~\cite{becker2015cia} introduced a
reliability-based CMA-ES attack that models XOR APUFs one component Arbiter PUF
at a time, breaking configurations up to 32-XOR (2M CRPs in simulation; 400k
CRPs on silicon-emulated 8-XOR-4-way data), and the CalyPSO
framework~\cite{calypso2024} combined particle swarm optimization with
hardware-friendly modeling, reporting attacks on XOR APUFs up to $k=20$.  While
these approaches established that higher XOR configurations are attackable, each
reported a single success point and none characterized the \emph{probability} of
success across the data spectrum.  Unlike prior work, we measure
$P(\text{success} \mid N, k)$ as a continuous function, revealing a probabilistic
rather than deterministic transition.

\paragraph{Deep learning attacks.}
Deep learning has become the dominant approach for XOR PUF modeling.
Santikellur and Chakraborty~\cite{santikellur2023} provided a comprehensive
survey and open-source implementation, establishing the standard 4-layer MLP
with Tanh hidden units as the de facto baseline.  Fei~et~al.~\cite{fei2024mope}
proposed the MoPE and MMoPE frameworks, which dynamically adjust model
complexity to match the PUF's XOR configuration, achieving success on up to
7-XOR with improved sample efficiency; their analysis emphasized that
architectural matching improves attack efficiency.  Sayadi~et~al.~\cite{sayadi2025}
introduced a chosen-challenge divide-and-conquer attack that models each
underlying Arbiter PUF individually using highly correlated challenge pairs,
breaking XOR APUFs even without reliability information.  Wang~et~al.~\cite{wang2024c2d2}
proposed the C2D2 framework, which exploits statistical covariances in the
challenge space and design dependency among PUF compositions to attack XOR
APUFs, iPUFs, and bistable ring PUFs on both simulation and FPGA.  While these
works advanced the state of the art in attack efficiency, they each frame the
problem as a search for the best possible architecture or training strategy.
Our work shows that below a data-dependent threshold, architectural
innovations provide no benefit, and above it, a standard MLP suffices.

\paragraph{Split attack and Interpose PUF.}
The split attack, introduced by Wisiol~et~al.~\cite{wisiol2020splitting}, exploits
structural rather than functional approximations: instead of learning the
complete $k$-XOR function, it decomposes it into two smaller
subproblems, training separate models on the upper and lower halves of the XOR
chain and combining their probabilistic outputs.  This reduces the effective
XOR complexity by approximately a factor of 2, drastically lowering the CRP
requirement.  Crucially, Wisiol~et~al. also demonstrated that the Interpose PUF
(iPUF)~\cite{nguyen2019}, proposed as a split-attack-resistant alternative,
remains vulnerable: a $(3,3)$-iPUF is attackable with 3M CRPs, refuting the
claimed security equivalence to a $(k_{\text{down}}+k_{\text{up}}/2)$-XOR
APUF.  We extend this result with
a systematic sweep showing that all configurations up to
$(4,4)$ fall at $10^6$ or fewer CRPs with $>99\%$ accuracy, consistent with the
$\max\{k_{\text{up}},k_{\text{down}}\}$ security level of the splitting attack,
moving beyond
the binary breakability statements typical of prior work.

\paragraph{Theoretical perspectives.}
The loss-landscape literature offers a theoretical perspective:
Choromanska~et~al.~\cite{choromanska2015loss}
showed that low-energy regions of deep-network loss landscapes form connected
manifolds, and Dauphin~et~al.~\cite{dauphin2014identifying}
identified saddle points as the dominant optimization obstruction in
high-dimensional nonconvex problems.  The grokking
phenomenon~\cite{power2022grokking}---where validation accuracy remains near
random for many epochs before suddenly rising to perfect
generalization---shares qualitative features with our data-driven
phase transition, but occurs along the epoch axis rather than the data axis.
A recent systematic review~\cite{pufmodeling2026systematic} of over 70
PUF modeling studies confirms that no prior work has characterized the
success probability as a function of CRP budget.  Our work connects these
theoretical insights to empirical PUF modeling, showing that the data-driven
phase transition arises from a growing attractor basin in the loss landscape,
captured quantitatively by the S-curve framework.


